\documentclass[12pt]{article}

\usepackage[a4paper, margin=2.5cm]{geometry}
\usepackage{amsmath}
\usepackage{amssymb}
\usepackage{amsthm}
\usepackage[colorlinks=true, linkcolor=blue, citecolor=blue, urlcolor=blue]{hyperref}
\usepackage{booktabs}
\usepackage{natbib}

\newtheorem{theorem}{Theorem}
\newtheorem{proposition}[theorem]{Proposition}
\newtheorem{corollary}[theorem]{Corollary}
\newtheorem{definition}[theorem]{Definition}
\newtheorem{remark}[theorem]{Remark}
\newtheorem{example}[theorem]{Example}
\newtheorem{lemma}[theorem]{Lemma}

\newcommand{\Ecal}{\mathcal{E}}
\newcommand{\Ocal}{\mathcal{O}}

\title{Observer Frames and Relativity\\from Event-Order Geometry}
\author{%
  Lee Smart\\[4pt]
  \textit{Institute of Vibrational Field Dynamics}\\[2pt]
  \texttt{contact@vibrationalfielddynamics.org}\\[2pt]
  \texttt{@vfd\_org}%
}
\date{April 2026}

\begin{document}

\maketitle

\noindent\textbf{Status:} Preprint (not peer-reviewed)

\begin{abstract}
Paper~XXIII derived a strict partial order on admissible overlap events between two phase-rotated 600-cell copies, interpreting the ordering as an emergent temporal structure. The present paper shows that this event poset carries observer-dependent structure analogous to two kinematic features of special relativity.

An observer is defined as a linear extension of the event partial order equipped with a positive weight function over events. Different observers assign different time functions to the same event set. The paper proves, for finite event posets admitting incomparable events: (1) relativity of simultaneity --- incomparable events can be ordered differently by different observers; (2) that sufficiently distinct observer weights yield unequal global clock intervals between the same comparable endpoints in a shared linear extension, providing a time-dilation analogue. It \emph{defines} (3) worldlines as maximal chains in the poset, and separately defines a chain-restricted time $\tau_O$ on them. A worked example on the XXIII dual-pentagon system exhibits all constructions explicitly.

The paper does not derive Lorentz invariance, the metric tensor, or the Einstein equations. It establishes that observer-dependent orderings and non-unique clock assignments exist on finite event posets with incomparable events, and exhibits the resulting relativity analogy concretely in the pentagon model of Paper~XXIII.
\end{abstract}

%==================================================================
\section{Introduction}\label{sec:intro}
%==================================================================

Paper~XXIII~\citep{SmartXXIII} established that the phase-overlap geometry of two conjugate 600-cell copies produces a strict partial order $\prec$ on an event set $\Ecal$. Paper~XXIII further shows, conditional on the existence of incomparable events (equivalently: at least two events sharing a birth angle), that no linear extension is canonical --- this is the no-canonical-clock conclusion. The conditional hypothesis is verified explicitly in the dual-pentagon toy model of Paper~XXIII but remains open for the full 600-cell event set. The present paper converts this structure into observer-dependent physics within the same conditional setting.

\medskip
\noindent\textbf{Standing assumption.} Throughout this paper, $\Ecal$ denotes a finite event set carrying a strict partial order $\prec$ (as in Paper~XXIII). All sums over events are therefore finite.
\medskip

The central question is: given a partial order on events with no preferred total ordering, how do different observers reconstruct temporal experience, and do the resulting differences reproduce features of relativistic kinematics?

\subsection*{Scope}

This paper derives observer-dependent time functions, proves relativity of simultaneity for incomparable event pairs, and proves a time-dilation analogue under a sharp hypothesis on differing weight sums. It does not derive Lorentz invariance, construct a metric, or claim equivalence with special relativity. It establishes observer-dependent orderings and non-unique clock assignments on finite event posets with incomparable events, and exhibits the resulting relativity analogy concretely in the dual-pentagon model of Paper~XXIII.

%==================================================================
\section{Observer Definition}\label{sec:observer}
%==================================================================

\begin{definition}[Observer]\label{def:observer}
An \textit{observer} on the event poset $(\Ecal, \prec)$ is a pair
\begin{equation}
    O = (\leq_O,\; \mu_O)
\end{equation}
where:
\begin{itemize}
    \item $\leq_O$ is a linear extension of $\prec$: a total order on $\Ecal$ such that $e_1 \prec e_2 \Rightarrow e_1 \leq_O e_2$,
    \item $\mu_O: \Ecal \to \mathbb{R}_{>0}$ is a \textit{sampling measure} (equivalently, a positive weight function): a positive function assigning a ``duration weight'' to each event.
\end{itemize}
\end{definition}

\begin{remark}
The linear extension $\leq_O$ determines \textit{what order} the observer assigns to events. The sampling measure $\mu_O$ determines \textit{how much time} the observer associates with each event. Together they produce a complete observer time function.
\end{remark}

\begin{remark}
No physical velocity, mass, or energy is assumed. The sampling measure is a structural degree of freedom intrinsic to the observer's relationship with the event poset.
\end{remark}

%==================================================================
\section{Observer Time Functions}\label{sec:time}
%==================================================================

\begin{definition}[Observer time]\label{def:time}
For an observer $O = (\leq_O, \mu_O)$, the \textit{observer time function} is:
\begin{equation}\label{eq:time}
    t_O(e) = \sum_{\substack{e' \in \Ecal \\ e' \leq_O e}} \mu_O(e')
\end{equation}
That is, the time assigned to event $e$ is the cumulative measure of all events up to and including $e$ in the observer's ordering.
\end{definition}

\begin{proposition}[Order preservation]\label{prop:preserve}
For any observer $O$:
\begin{equation}
    e_1 \prec e_2 \;\Rightarrow\; t_O(e_1) < t_O(e_2)
\end{equation}
\end{proposition}

\begin{proof}
If $e_1 \prec e_2$, then $e_1 \leq_O e_2$ (since $\leq_O$ extends $\prec$) and $e_1 \neq e_2$. The sum defining $t_O(e_2)$ includes all terms in $t_O(e_1)$ plus at least $\mu_O(e_2) > 0$. Therefore $t_O(e_2) > t_O(e_1)$.
\end{proof}

\begin{remark}
The time function $t_O$ is not unique: it depends on both the choice of linear extension and the choice of sampling measure. Different observers generically assign different times to the same events, even when they agree on the partial order.
\end{remark}

%==================================================================
\section{Relativity of Simultaneity}\label{sec:simultaneity}
%==================================================================

\begin{theorem}[Relativity of simultaneity]\label{thm:simultaneity}
If events $e_1, e_2 \in \Ecal$ are incomparable under $\prec$ (neither $e_1 \prec e_2$ nor $e_2 \prec e_1$), then there exist observers $O_1, O_2$ such that:
\begin{equation}
    t_{O_1}(e_1) < t_{O_1}(e_2) \qquad\text{and}\qquad t_{O_2}(e_2) < t_{O_2}(e_1)
\end{equation}
\end{theorem}

\begin{proof}
Since $e_1$ and $e_2$ are incomparable, form two augmented relations $\prec_1 := \prec \cup \{(e_1,e_2)\}$ and $\prec_2 := \prec \cup \{(e_2,e_1)\}$. Each remains irreflexive and acyclic: if, e.g., $\prec_1$ admitted a cycle it would force $e_2 \prec e_1$, contradicting incomparability. Taking transitive closures yields strict partial orders extending $\prec$, and Szpilrajn's order-extension theorem~\citep{Szpilrajn1930} gives linear extensions $\leq_1,\leq_2$ of these, placing $e_1$ before $e_2$ and $e_2$ before $e_1$ respectively (cf.\ Paper~XXIII's no-canonical-clock theorem~\citep{SmartXXIII}, which uses exactly this transitive-closure-plus-Szpilrajn construction for the same event set). Choose any positive weight function $\mu > 0$ and set $O_1 = (\leq_1, \mu)$, $O_2 = (\leq_2, \mu)$. Since $\leq_1$ places $e_1$ strictly before $e_2$, the sum in~\eqref{eq:time} defining $t_{O_1}(e_2)$ includes every term in $t_{O_1}(e_1)$ plus at least $\mu(e_2) > 0$, so $t_{O_1}(e_1) < t_{O_1}(e_2)$; the same argument with $\leq_2$ gives $t_{O_2}(e_2) < t_{O_2}(e_1)$.
\end{proof}

\begin{remark}
This is the geometric analogue of the relativity of simultaneity in special relativity: events that are spacelike-separated (incomparable in the causal order) can be ordered differently by different inertial frames (linear extensions). The VFD framework derives this from the event poset; special relativity derives it from Lorentz invariance. The structural role is analogous.
\end{remark}

%==================================================================
\section{Time Dilation Analogue}\label{sec:dilation}
%==================================================================

\begin{definition}[Observer endpoint interval of a chain]\label{def:elapsed}
For a chain $C = (e_1 \prec e_2 \prec \cdots \prec e_n) \subset \Ecal$ and observer $O = (\leq_O, \mu_O)$, the \textit{observer endpoint interval} is the global clock difference at the endpoints of $C$:
\begin{equation}\label{eq:elapsed}
    \Delta t_O(C) \;=\; t_O(e_n) - t_O(e_1) \;=\; \sum_{\substack{e' \in \Ecal \\ e_1 <_O e' \,\leq_O\, e_n}} \mu_O(e').
\end{equation}
Only the endpoints $e_1, e_n$ of $C$ enter this definition; the intermediate events of the chain play no role, and the sum ranges over all events in the observer's total-order interval $(e_1, e_n]$.
\end{definition}

\begin{theorem}[Observer-dependent endpoint interval under unequal weight sums]\label{thm:dilation}
Let $\leq$ be a linear extension of $\prec$, and let $O_1 = (\leq, \mu_1)$, $O_2 = (\leq, \mu_2)$ be observers sharing this extension but with different sampling measures. Let $C = (e_1 \prec \cdots \prec e_n)$ be a chain, and let
\[
    S \;=\; \{\,e' \in \Ecal : e_1 < e' \leq e_n\,\}
\]
(using the shared order $\leq$, with strict part $<$) be the set of events strictly after $e_1$ and up to $e_n$. If
\begin{equation}\label{eq:dilation-hyp}
    \sum_{e' \in S} \mu_1(e') \;\neq\; \sum_{e' \in S} \mu_2(e'),
\end{equation}
then the global clock intervals differ:
\begin{equation}
    \Delta t_{O_1}(C) \neq \Delta t_{O_2}(C).
\end{equation}
Both observers agree on the ordering of events; they disagree on the cumulative clock reading between the endpoints of the chain.
\end{theorem}

\begin{proof}
Since $O_1$ and $O_2$ share the linear extension $\leq$, the set $S$ is identical for both. By Definition~\ref{def:elapsed},
\[
    \Delta t_{O_i}(C) = \sum_{e' \in S} \mu_i(e') \qquad (i=1,2).
\]
Hypothesis~\eqref{eq:dilation-hyp} gives the conclusion directly.
\end{proof}

\begin{remark}
The hypothesis~\eqref{eq:dilation-hyp} is an inequality of scalar sums, not a pointwise statement; cancellations between individual-event weight differences are allowed so long as the sums do not coincide. A simple sufficient condition is pointwise dominance $\mu_1(e') \geq \mu_2(e')$ on $S$ with strict inequality on at least one event.
\end{remark}

\begin{remark}
This is the event-order analogue of time dilation at the level of the global clock: two observers sharing a linear extension of the event poset, but with different sampling measures, assign different cumulative clock intervals between the same pair of comparable endpoints. In special relativity, frame-dependent time intervals between the same pair of spacetime events arise from relative velocity; here they arise from the difference in weight sums~\eqref{eq:dilation-hyp}. The analogue concerns global clock readings (not intrinsic chain/proper time $\tau_O$, which is defined separately in Section~\ref{sec:worldlines}).
\end{remark}

\begin{remark}
The sampling measure $\mu_O$ is not yet connected to physical velocity or energy. Establishing that connection requires further structure (the metric emergence of Paper~XXV is a candidate step, but no specific result of XXV is used here). The present paper demonstrates only that the \textit{kinematic possibility} of a time-dilation analogue exists in the event-order framework.
\end{remark}

%==================================================================
\section{Worldlines}\label{sec:worldlines}
%==================================================================

\begin{definition}[Worldline]\label{def:worldline}
A \textit{worldline} is a maximal chain in $(\Ecal, \prec)$: a totally ordered subset $W \subset \Ecal$ that cannot be extended by inclusion.
\end{definition}

\begin{remark}
A worldline is a complete sequence of events each following the previous in the partial order. Different worldlines may share some events (intersect) or be entirely disjoint.
\end{remark}

\begin{definition}[Chain time]\label{def:chaintime}
For an observer $O$ and a chain $W = (e_1 \prec \cdots \prec e_n)$, the \textit{chain time} (or worldline time) is:
\begin{equation}
    \tau_O(W) = \sum_{k=1}^{n} \mu_O(e_k)
\end{equation}
This is the cumulative measure restricted to events on the chain itself.
\end{definition}

\begin{remark}
Chain time $\tau_O(W)$ and observer elapsed time $\Delta t_O(W) = t_O(e_n) - t_O(e_1)$ are distinct cumulative-weight constructions: $\tau_O$ sums weights along the chain $W$ itself (including its initial event), while $\Delta t_O$ sums weights over all events strictly after $e_1$ and up to $e_n$ in the observer's total ordering, irrespective of whether they lie on $W$. Both depend on the observer $O$; neither is shown in this paper to be observer-independent or intrinsic to $W$.
\end{remark}

%==================================================================
\section{Relative Velocity (Proto-Concept)}\label{sec:velocity}
%==================================================================

A full notion of velocity requires spatial structure, which this paper does not construct. However, a proto-concept can be defined.

\begin{definition}[Relative sampling rate]\label{def:velocity}
For two observers $O_1, O_2$, their \textit{relative sampling rate} along a chain $C$ is:
\begin{equation}
    \rho(O_1, O_2; C) = \frac{\tau_{O_1}(C)}{\tau_{O_2}(C)} = \frac{\sum_{e \in C} \mu_1(e)}{\sum_{e \in C} \mu_2(e)}
\end{equation}
\end{definition}

When $\rho \neq 1$, the two observers assign different chain-restricted cumulative weights to the same chain. This is a structural precursor to relative velocity: a dimensionless comparison of chain-weight totals between observers.

\begin{remark}
This proto-velocity does not have units of distance/time. It is a dimensionless ratio of sampling measures. Connecting it to physical velocity requires the spatial metric structure of Paper~XXV.
\end{remark}

%==================================================================
\section{Worked Example: Dual Pentagons}\label{sec:example}
%==================================================================

\begin{example}[Extending the pentagon system]\label{ex:pentagon}
We use the dual-pentagon event set of Paper~XXIII~\citep{SmartXXIII}: its first two birth-angle layers
\[
    L_1 \;=\; \{(k, k-1 \bmod 5) : k=0,\dots,4\} \quad \text{at } \theta_b = \pi/5 - \delta,
\]
\[
    L_2 \;=\; \{(k, k-2 \bmod 5) : k=0,\dots,4\} \quad \text{at } \theta_b = 3\pi/5 - \delta,
\]
i.e.\ $L_1 = \{(0,4),(1,0),(2,1),(3,2),(4,3)\}$ and $L_2 = \{(0,3),(1,4),(2,0),(3,1),(4,2)\}$. The partial order $\prec$ compares birth angles, so every event in $L_1$ strictly precedes every event in $L_2$; events within a single layer are incomparable.

\textbf{Observers.} Fix the linear extension
\[
    \leq_A:\quad (0,4) < (1,0) < (2,1) < (3,2) < (4,3) < (0,3) < (1,4) < (2,0) < (3,1) < (4,2),
\]
and define three observers:
\begin{itemize}
    \item $O_1 = (\leq_A, \mu_1)$ with $\mu_1 \equiv 1$ (uniform).
    \item $O_2 = (\leq_B, \mu_2)$ with $\leq_B$ obtained from $\leq_A$ by swapping the incomparable pair $(0,4)$ and $(1,0)$, and $\mu_2 \equiv 1$.
    \item $O_3 = (\leq_A, \mu_3)$ with $\mu_3(e) = 2$ for $e \in L_1$ and $\mu_3(e) = 1$ for $e \in L_2$.
\end{itemize}
All three orderings respect $\prec$ (layer 1 before layer 2).

\textbf{Simultaneity (compare $O_1$ vs $O_2$).} The events $(0,4)$ and $(1,0)$ are incomparable. Observer $O_1$ assigns $t_{O_1}(0,4) = 1 < t_{O_1}(1,0) = 2$. Observer $O_2$ assigns $t_{O_2}(1,0) = 1 < t_{O_2}(0,4) = 2$. The same incomparable pair is ordered oppositely by the two observers; this instantiates Theorem~\ref{thm:simultaneity}.

\textbf{Time dilation (compare $O_1$ vs $O_3$).} $O_1$ and $O_3$ share the linear extension $\leq_A$. Consider the chain $C = ((0,4), (0,3))$, valid because $(0,4) \in L_1$ precedes $(0,3) \in L_2$. The set $S$ of events strictly after $(0,4)$ and up to $(0,3)$ in $\leq_A$ is
\[
    S = \{(1,0),(2,1),(3,2),(4,3),(0,3)\}.
\]
Hypothesis~\eqref{eq:dilation-hyp} holds: $\sum_S \mu_1 = 5$, $\sum_S \mu_3 = 2+2+2+2+1 = 9$. So
\[
    \Delta t_{O_1}(C) = 5, \qquad \Delta t_{O_3}(C) = 9,
\]
instantiating Theorem~\ref{thm:dilation}. Both observers agree on the ordering; they disagree on elapsed time.

\textbf{Relative sampling rate along $C$.} By Definition~\ref{def:velocity}, the rate uses the chain-restricted time $\tau$:
\[
    \tau_{O_1}(C) = \mu_1(0,4) + \mu_1(0,3) = 1 + 1 = 2, \qquad
    \tau_{O_3}(C) = \mu_3(0,4) + \mu_3(0,3) = 2 + 1 = 3,
\]
hence $\rho(O_1, O_3; C) = 2/3$. The chain-time ratio differs from the global elapsed-time ratio $\Delta t_{O_1}/\Delta t_{O_3} = 5/9$; the two are distinct quantities built from different sums (chain events vs.\ all ordering-interval events), and in general there is no algebraic relation between them.
\end{example}

%==================================================================
\section{Relation to Special Relativity}\label{sec:SR}
%==================================================================

\begin{table}[ht]
\centering
\caption{Structural comparison: VFD event-order framework vs.\ special relativity.}
\label{tab:SR}
\begin{tabular}{@{}lll@{}}
\toprule
Feature & VFD (this paper) & Special relativity \\
\midrule
Events & Admissible vertex pairs & Spacetime events \\
Ordering & Birth-angle partial order & Causal structure \\
Observer & Linear extension + weight function & Inertial frame \\
Global clock $t_O(e)$ & Cumulative weights over $\leq_O$ & Coordinate time \\
Chain time $\tau_O(W)$ & Weights along the chain only & Proper-time analogue \\
Simultaneity & Observer-dependent (Thm.~\ref{thm:simultaneity}) & Frame-dependent \\
Time dilation & Weight-sum differences (Thm.~\ref{thm:dilation}) & Velocity-dependent \\
Worldline & Maximal chain in poset (defined) & Timelike curve \\
Velocity & Proto-concept (chain-time ratio $\rho$) & $dx/dt$ \\
\bottomrule
\end{tabular}
\end{table}

\begin{remark}[Epistemic status]
The structural analogies in Table~\ref{tab:SR} are \textit{exhibited}, not derived from Lorentz invariance. The VFD framework produces two qualitative phenomena (observer-dependent simultaneity for incomparable event pairs, and observer-dependent endpoint intervals under unequal weight sums) from a different starting point (event poset with sampling measures), and defines a notion of worldlines as maximal chains. Whether this structural analogy can be promoted to full Lorentz covariance requires metric structure not yet available.
\end{remark}

%==================================================================
\section{Discussion}\label{sec:discussion}
%==================================================================

\subsection{What is established}

\begin{itemize}
    \item Observer defined as linear extension + sampling measure (Definition~\ref{def:observer}).
    \item Observer time function constructed explicitly (Definition~\ref{def:time}, Proposition~\ref{prop:preserve}).
    \item Relativity of simultaneity proved for incomparable event pairs (Theorem~\ref{thm:simultaneity}).
    \item Time-dilation analogue proved under the sharp weight-sum hypothesis~\eqref{eq:dilation-hyp} (Theorem~\ref{thm:dilation}).
    \item Worldlines \emph{defined} as maximal chains (Definition~\ref{def:worldline}); not a theorem.
    \item Worked pentagon example exhibits all constructions with fully checked arithmetic.
\end{itemize}

\subsection{What is not established}

\begin{itemize}
    \item No Lorentz invariance. The framework has discrete ($H_4$) symmetry, not continuous Lorentz symmetry.
    \item No metric. The sampling measure gives duration differences but not distances.
    \item No Einstein equations. The connection between observer structure and gravitational dynamics is not addressed.
    \item No physical velocity. The relative sampling rate is a structural precursor, not a derived physical quantity.
    \item The sampling measure $\mu_O$ is a structural degree of freedom, not yet connected to physical observables (mass, energy, motion).
\end{itemize}

\subsection{Programme position}

\begin{itemize}
    \item Paper~XXIII: event structure and emergent temporal ordering.
    \item Paper~XXIV (this paper): observer frames and relativity-like kinematics.
    \item Paper~XXV (future): metric structure and the route to dynamics.
\end{itemize}

%==================================================================
\section{Conclusion}\label{sec:conclusion}
%==================================================================

In the pentagon model of Paper~XXIII, and more generally for any finite event poset admitting incomparable events, the event poset equipped with observer-dependent weight functions supports a kinematic precursor of relativistic structure: observer-dependent simultaneity for incomparable event pairs, and an observer-dependent endpoint interval between the same pair of comparable endpoints under unequal weight sums. These two facts are proved under sharp, explicit hypotheses (Theorems~\ref{thm:simultaneity} and~\ref{thm:dilation}) and exhibited concretely on the pentagon model. A notion of worldlines is introduced by defining them as maximal chains, with chain-restricted cumulative weight $\tau_O$ as a separately defined quantity. No spacetime manifold, metric, or Lorentz symmetry is postulated, and no claim is made about incomparable events in the full 600-cell event set (which remains open in Paper~XXIII).

The key structural insight is that the non-uniqueness of a global clock in models with incomparable events (Paper~XXIII, pentagon) is not a deficiency but a feature: it is a geometric source of observer dependence in this framework. Different observers are different ways of completing the partial order into a total order, and different weight functions produce different assigned clock intervals between the same comparable endpoints.

The next step is to equip the event structure with spatial and metric content (Paper~XXV), connecting the sampling measure to physical observables and determining whether the full Lorentz-covariant structure can be recovered from the 600-cell geometry.

%==================================================================
\bibliographystyle{plainnat}
\bibliography{references}

\end{document}
